\documentclass[10pt,a4paper]{article}
% Preâmbulo comum: Operação Sentinel (FEELT37802 - Sistemas Distribuídos 2026/2)
\PassOptionsToPackage{hyphens,spaces,obeyspaces}{url}
\usepackage[utf8]{inputenc}
\usepackage[T1]{fontenc}
\usepackage{lmodern}
\usepackage[scaled=0.95]{helvet}   % fonte sem serifa (Helvetica/Arial) no documento todo
\renewcommand{\familydefault}{\sfdefault}
\usepackage[brazilian]{babel}
\usepackage[a4paper,margin=2cm,top=2.2cm,bottom=2cm]{geometry}
\usepackage{needspace}
\usepackage{microtype}
\usepackage{amssymb}
\usepackage{xcolor}
\usepackage{graphicx}
\usepackage{array}
\usepackage{tabularx}
\usepackage{longtable}
\usepackage{booktabs}
\usepackage{enumitem}
\usepackage{fancyhdr}
\usepackage{listings}
\usepackage[most]{tcolorbox}
\usepackage{tikz}
\usetikzlibrary{arrows.meta,positioning,calc,fit,shapes.misc,decorations.pathmorphing}
\usepackage[colorlinks=true,linkcolor=blue!50!black,urlcolor=blue!50!black,citecolor=blue!50!black]{hyperref}

\setlength{\parindent}{0pt}
\setlength{\parskip}{3pt}
\setlist{itemsep=1pt,topsep=3pt,parsep=0pt}
\renewcommand{\arraystretch}{1.1}

% ---- cores ----
\definecolor{azul}{HTML}{1F4E79}
\definecolor{verde}{HTML}{2E7D32}
\definecolor{laranja}{HTML}{C75B00}
\definecolor{vermelho}{HTML}{B3261E}
\definecolor{cinza}{HTML}{F3F4F6}
\definecolor{cinzaescuro}{HTML}{555555}

% ---- comandos utilitários ----
\usepackage{url}
\DeclareUrlCommand\cmd{\urlstyle{tt}}
\makeatletter
\def\UrlBreaks{\do\/\do\-\do\_\do\.\do\=\do\:\do\,\do\)\do\(}
\makeatother
\Urlmuskip=0mu plus 2mu
\newcommand{\seta}{$\rightarrow$}
\newcommand{\caixa}{\raisebox{-0.15ex}{$\square$}}
\newcommand{\linhas}[1]{\par\vspace{2pt}\foreach \i in {1,...,#1}{\noindent\rule{\linewidth}{0.3pt}\\[1.15em]}}

% ---- cabeçalho/rodapé (estilo: código e disciplina em negrito à esquerda, professor à direita) ----
\pagestyle{fancy}
\fancyhf{}
\renewcommand{\headrulewidth}{0.4pt}
\fancyhead[L]{\footnotesize\textbf{FEELT37802 - SISTEMAS DISTRIBUÍDOS (2026/2)}}
\fancyhead[R]{\footnotesize Prof. Renato Carrijo}
\fancyfoot[C]{\small\thepage}

% ---- bloco de título: centralizado, negrito, subtítulo menor e filete ----
\newcommand{\titulosentinel}[2]{%
  \begin{center}
    {\LARGE\bfseries #1\par}\vspace{6pt}
    {\small #2\par}
  \end{center}\vspace{-6pt}
  \noindent\rule{\linewidth}{0.5pt}\par\vspace{8pt}}

% ---- bloco da equipe ----
\newcommand{\blocoequipe}{%
  \begin{tcolorbox}[colback=cinza,colframe=black!30,boxrule=0.5pt,arc=2pt,left=8pt,right=8pt,top=6pt,bottom=6pt]
  \noindent\textbf{Equipe:}\ \rule[-2pt]{12.6cm}{0.4pt}\par\vspace{7pt}
  \noindent\textbf{Líder (quem entrega):}\ \rule[-2pt]{11.1cm}{0.4pt}\par\vspace{7pt}
  \begin{tabular}{@{}>{\centering\arraybackslash}p{0.8cm}@{\hspace{4pt}}p{8.4cm}@{\hspace{8pt}}p{6.2cm}@{}}
  \textbf{N.} & \textbf{Integrante} & \textbf{Papel/Função}\\[3pt]
  1 & \dotfill & \dotfill\\[5pt]
  2 & \dotfill & \dotfill\\[5pt]
  3 & \dotfill & \dotfill\\[5pt]
  4 & \dotfill & \dotfill
  \end{tabular}
  \end{tcolorbox}}

% ---- listagens ----
\lstdefinestyle{shell}{
  basicstyle=\ttfamily\small,
  backgroundcolor=\color{cinza},
  frame=single, rulecolor=\color{black!25},
  columns=fullflexible, keepspaces=true, breaklines=true, breakatwhitespace=false,
  postbreak=\mbox{\textcolor{cinzaescuro}{$\hookrightarrow$}\space},
  xleftmargin=2pt, xrightmargin=2pt, aboveskip=6pt, belowskip=4pt,
  showstringspaces=false, upquote=true,
  commentstyle=\color{verde}\itshape,
  morecomment=[l]{\#},
  literate=
    {á}{{\'a}}1 {é}{{\'e}}1 {í}{{\'i}}1 {ó}{{\'o}}1 {ú}{{\'u}}1
    {Á}{{\'A}}1 {É}{{\'E}}1 {Í}{{\'I}}1 {Ó}{{\'O}}1 {Ú}{{\'U}}1
    {â}{{\^a}}1 {ê}{{\^e}}1 {ô}{{\^o}}1 {Â}{{\^A}}1 {Ê}{{\^E}}1 {Ô}{{\^O}}1
    {ã}{{\~a}}1 {õ}{{\~o}}1 {Ã}{{\~A}}1 {Õ}{{\~O}}1
    {à}{{\`a}}1 {ç}{{\c c}}1 {Ç}{{\c C}}1 {ü}{{\"u}}1
}
\lstset{style=shell}
\lstnewenvironment{terminal}{\lstset{style=shell}}{}
\lstdefinestyle{saida}{style=shell, backgroundcolor=\color{white}, frame=leftline, morecomment=[l]{\#}}
\lstnewenvironment{saida}{\lstset{style=saida}}{}

% ---- caixas ----
\newtcolorbox{ideia}[1][Ideia-chave]{enhanced,breakable,colback=azul!6,colframe=azul,
  fonttitle=\bfseries,title={#1},boxrule=0.6pt,arc=2pt,left=6pt,right=6pt,top=3pt,bottom=3pt}
\newtcolorbox{atencao}[1][Atenção]{enhanced,breakable,colback=laranja!7,colframe=laranja,
  fonttitle=\bfseries,title={#1},boxrule=0.6pt,arc=2pt,left=6pt,right=6pt,top=3pt,bottom=3pt}
\newtcolorbox{perigo}[1][Cuidado]{enhanced,breakable,colback=vermelho!6,colframe=vermelho,
  fonttitle=\bfseries,title={#1},boxrule=0.6pt,arc=2pt,left=6pt,right=6pt,top=3pt,bottom=3pt}
\newtcolorbox{prever}[1][Preveja antes de executar]{enhanced,breakable,colback=verde!6,colframe=verde,
  fonttitle=\bfseries,title={#1},boxrule=0.6pt,arc=2pt,left=6pt,right=6pt,top=3pt,bottom=3pt}
\newtcolorbox{evid}[1][Evidência a registrar]{enhanced,breakable,colback=white,colframe=black!55,
  fonttitle=\bfseries,title={#1},boxrule=0.6pt,arc=2pt,left=6pt,right=6pt,top=3pt,bottom=3pt}
\newtcolorbox{checagem}[1][Pergunta de checagem]{enhanced,breakable,colback=black!3,colframe=black!40,
  fonttitle=\bfseries,title={#1},boxrule=0.5pt,arc=2pt,left=6pt,right=6pt,top=3pt,bottom=3pt}

% ---- títulos de hora ----
\newcommand{\hora}[3]{%
  \Needspace{14\baselineskip}\vspace{6pt}\phantomsection
  \begin{tcolorbox}[enhanced,colback=azul,colframe=azul,coltext=white,arc=3pt,boxrule=0pt,left=8pt,right=8pt,top=5pt,bottom=5pt]
    {\Large\bfseries Hora #1 -- #2}\hfill {\normalsize #3}
  \end{tcolorbox}
  \addcontentsline{toc}{section}{Hora #1 -- #2}
  \vspace{2pt}}

\begin{document}

\titulosentinel{Relatório de Incidente -- Operação Sentinel (Entrega F)}{FEELT37802 - Sistemas Distribuídos (2026/2) \,·\, Prof. Renato Carrijo \,·\, Equipe: Heitor Carvalho Andrade (líder, 12311ECP023) e integrantes}

\textbf{Incidente:} ``O servidor está vivo, mas o sistema não confia nele'' (Hora 7), conduzido com \cmd{scripts/inc.sh iniciar} em
\textbf{19/11/2026, 03:27 a 03:31 UTC}, com tráfego de fundo de 3~req/s. Nesta execução o nó afetado foi o \textbf{Rust~B} (o enunciado
descreve o roteiro com o Rust~A; os eventos são os mesmos, trocando o nó). Todos os horários são UTC e vêm dos arquivos
\cmd{evidencias/hora7-*}, \cmd{evidencias/20261119-032645_linha-de-base/} e \cmd{evidencias/20261119-032929_antes-do-reinicio/}.

\section{Resumo em três linhas}
O watchdog marcou o Rust~B como DOWN às 03:27:35 (três sondagens \cmd{/health/ready} com HTTP~503) enquanto o Nginx seguia entregando 200 por B;
aos 03:28:42 surgiram retries e latências de 1~s, causadas por uma regra de rede (atraso 150~ms, perda 25\%) na saída de B para o Nginx.
A equipe preservou evidências, decidiu \emph{retirar B do balanceador} (ação reversível) em vez de reiniciar, localizou as duas causas imediatas
(prontidão desligada em B e regra de rede) e confirmou a recuperação às 03:31:00 sem nenhum reinício de contêiner.

\section{Linha do tempo}
\begin{center}\small
\begin{tabularx}{\linewidth}{@{}>{\ttfamily}l>{\raggedright\arraybackslash}X>{\raggedright\arraybackslash}p{4.6cm}@{}}
\toprule
\textbf{UTC} & \textbf{O que aconteceu} & \textbf{Evidência}\\ \midrule
03:26:35 & Estado normal confirmado: 7 contêineres \cmd{healthy}, watchdog A e B UP, nenhum alerta, nenhuma regra de rede (\cmd{chaos.sh status todos}). Linha de base congelada. & \cmd{hora7-00-antes.txt}; \cmd{20261119-032645_linha-de-base/}\\
03:27:00 & Início do exercício (\cmd{inc.sh iniciar}). & \cmd{hora7-eventos-inc.txt}\\
03:27:25--35 & Watchdog: três sondagens seguidas de B com \cmd{reason=http_503}; \cmd{state_change rust-b up\seta{}down} às 03:27:35,075. & \cmd{log-watchdog.txt} (antes-do-reinicio)\\
03:27:41 & \cmd{NoDown\{rust-b\}} entra em \emph{pending}; dispara às 03:27:56 (\cmd{for: 15s}). & \cmd{hora7-alertas-prometheus.txt}\\
03:27:47 & Painel: ``rust-b está DOWN segundo o watchdog: 4 falhas seguidas (motivo: http\_503)''. & \cmd{hora7-painel-20261119-032747Z-rust-b-down.jpg}\\
03:28:02 & \textbf{Checagem 1.} Usuário (60~s): 176 req, 100\% 200, 50\%/50\% A/B, 0 retry, p95 1~ms. De um \emph{terceiro ponto} (toolbox): \cmd{rust-b/health/ready \seta{} 503}, \cmd{rust-b/health/live \seta{} 200}; \cmd{admin status}: \cmd{rust-b ready_flag=false effective_ready=false delay_ms=0 uptime_s=1441}. Idade da última sondagem $<5$~s; \cmd{up} = 1 nos 4 alvos. & \cmd{hora7-01-checagem-1.txt}\\
03:28:10 & Evento anunciado: ``o Nginx continua registrando 200 em requisições encaminhadas a B'' (confirmado: linhas \cmd{node=rust-b status=200} às 03:27:46). & \cmd{hora7-eventos-inc.txt}; \cmd{hora7-01-checagem-1.txt}\\
03:28:37 & \textbf{Checagem 2.} Quadro igual ao da checagem 1 (sem latência, sem retry). & \cmd{hora7-02-checagem-2.txt}\\
03:28:42 & Evento anunciado: ``o tempo de resposta aumenta; requisições repetidas em outro nó''. & \cmd{hora7-eventos-inc.txt}\\
03:29:11 & \textbf{Checagem 3.} Usuário (60~s): 5 retries (2,8\%), máx. 1,005~s, B caiu para 35\% do tráfego; \cmd{chaos.sh status}: em \cmd{rust-b}, \cmd{[saida -> nginx (10.9.8.10)] delay 150ms loss 25\%}; nos demais contêineres, nenhuma regra. Watchdog: B DOWN, \cmd{falhas=19}, ainda \cmd{http_503} (sondagens respondidas em 4~ms: o caminho watchdog\seta{}B não tem a regra). & \cmd{hora7-03-checagem-3.txt}\\
03:29:13--23 & Captura na interface de B: SYN do Nginx às 13,616 respondido às 14,782 (1,17~s) e recusado com RST (o Nginx desistiu em 1~s); \emph{handshakes} com $\sim$150~ms; sondagem do watchdog (10.9.8.20) respondida de imediato. & \cmd{hora7-04-captura-rust-b.txt}\\
03:29:29 & Regra de ouro: evidências congeladas \textbf{antes} de qualquer ação (\cmd{evidence.sh antes-do-reinicio --rede}). & \cmd{20261119-032929_antes-do-reinicio/}\\
03:29:30 & Evento anunciado: ``o painel mostra B em vermelho, mas o Nginx segue incluindo B''. & \cmd{hora7-eventos-inc.txt}\\
03:29:44 & \textbf{Decisão:} drenar B do Nginx (\cmd{admin.sh pool drenar rust-b}), sem reiniciar nada. O comando \textbf{não surtiu efeito} nesta máquina (macOS: o \cmd{sed -i} do script é GNU); detectado na hora pelo \cmd{pool status} ``depois'' idêntico ao ``antes''. & \cmd{hora7-05-decisao-drenar-b.txt}\\
03:30:00 & Fim da injeção (o sistema fica como está até alguém agir). & \cmd{hora7-eventos-inc.txt}\\
03:30:17 & \textbf{Checagem após a decisão:} B ainda com 15\% do tráfego, 9 retries, p95 1,002~s, máx. 2,014~s: confirma que a drenagem não ocorreu. Causas imediatas já identificadas nas checagens; \textbf{correção dirigida}: \cmd{admin.sh ready rust-b 1} e \cmd{chaos.sh limpar rust-b}. & \cmd{hora7-06-apos-drenagem.txt}; \cmd{hora7-07-correcao.txt}\\
03:30:25 & Watchdog: \cmd{state_change rust-b down\seta{}up} (duas sondagens boas). 03:30:27 \cmd{pool restaurar rust-b} (sem efeito, pois B nunca saiu). 03:30:32: nenhum alerta ativo. & \cmd{log-watchdog}; \cmd{hora7-08-restaurar-b.txt}; \cmd{hora7-alertas-prometheus.txt}\\
03:31:00 & \textbf{Recuperação confirmada}: 20 req \seta{} \cmd{rust-a=9 rust-b=11}, 100\% 200; usuário (30~s): 108 req, 0 retry, p95 1~ms; \cmd{reinicios=0} em nginx, rust-a, rust-b e watchdog; nenhuma regra de rede. 03:31:11: \cmd{inc.sh resolver}. & \cmd{hora7-09-recuperacao-confirmada.txt}\\
03:33:10 & Pós-incidente: \cmd{scripts/admin.sh} corrigido (sed portável) e a ação reversível \textbf{provada}: com B drenado, 20/20 requisições por A; restaurado, 10/10. & \cmd{hora7-10-teste-drenagem-reversivel.txt}\\
\bottomrule
\end{tabularx}
\end{center}

\section{Fatos versus inferências}
\begin{center}\small
\begin{tabularx}{\linewidth}{@{}>{\raggedright\arraybackslash}X>{\raggedright\arraybackslash}X@{}}
\toprule
\textbf{Fatos (com evidência e horário)} & \textbf{Inferências (o que \emph{não} se podia afirmar)}\\ \midrule
O watchdog recebeu \textbf{HTTP 503} de \cmd{rust-b:8080/health/ready} em três sondagens seguidas e marcou DOWN às 03:27:35 (\cmd{reason=http_503}, não \cmd{timeout}). &
``O Rust~B caiu/travou'': \textbf{falso}. \cmd{/health/live} 200, \cmd{uptime_s} crescendo, Nginx entregando 200 por B, \cmd{docker compose ps} \cmd{healthy}.\\ \midrule
O Nginx continuou encaminhando a B e recebendo 200 (\cmd{node=rust-b status=200} às 03:27:46; 50\%/50\% às 03:28:02). &
``O watchdog está errado / com falha de rede'': \textbf{falso}. Um terceiro ponto (toolbox) viu o mesmo 503, e o próprio nó informou \cmd{ready_flag=false}.\\ \midrule
De um terceiro ponto: \cmd{ready \seta{} 503} e \cmd{live \seta{} 200}; \cmd{admin status}: \cmd{ready_flag=false} em B (03:28:02). &
``O monitor parou'' (evento 4 do enunciado): \textbf{não ocorreu nesta execução}: \cmd{up\{job="watchdog"\}=1}, idade da sondagem $<5$~s, painel atualizando.\\ \midrule
A partir de 03:28:42: retries \cmd{upstream_status=504, 200} com \cmd{req_time≈1,0 s}; \cmd{chaos.sh status} mostrou \cmd{delay 150ms loss 25\%} na saída de B \textbf{só para o Nginx}; a captura mostrou SYN-ACK atrasado 1,17~s e RST do Nginx. &
``A latência vem do nó B sobrecarregado'': \textbf{falso}. \cmd{delay_ms=0}, p95 \emph{no nó} em milissegundos, e a sondagem do watchdog para B respondia em 4~ms (caminho sem a regra).\\ \midrule
A regra de rede \textbf{não constava} em \cmd{evidencias/chaos.log} (nossos scripts registram toda injeção): foi aplicada por outro agente. &
``Alguém da equipe rodou \cmd{chaos.sh}'': refutado pelo log; a origem exata (o exercício) está fora do que o sistema registra.\\
\bottomrule
\end{tabularx}
\end{center}

\section{Hipóteses e como foram testadas}
\begin{enumerate}[leftmargin=*]
  \item \textbf{B marcou-se ``não pronto'' (prontidão desligada), processo vivo.} Teste: \cmd{curl} de um terceiro ponto a \cmd{/health/ready} e \cmd{/health/live} e \cmd{admin.sh status}. Resultado: 503/200 e \cmd{ready_flag=false}. \textbf{Confirmada} (03:28:02).
  \item \textbf{Partição ou perda de rede entre watchdog e B} (o observador erra, o nó está bem). Teste: a mesma sondagem a partir da toolbox; o motivo registrado pelo watchdog. Resultado: toolbox viu o mesmo 503; motivo \cmd{http_503} em 1--4~ms (há conectividade e resposta HTTP). \textbf{Refutada}.
  \item \textbf{B morreu ou reiniciou em loop.} Teste: \cmd{live}, \cmd{uptime_s}, \cmd{docker compose ps}, \cmd{RestartCount}. Resultado: \cmd{live} 200, \cmd{uptime_s=1441} e crescendo, \cmd{healthy}, \cmd{reinicios=0}. \textbf{Refutada}.
  \item \textbf{Degradação de rede no caminho Nginx$\leftrightarrow$B} (surgiu com o evento das 03:28:42). Teste: \cmd{status.sh --usuario} (retries, \cmd{req_time}), \cmd{chaos.sh status todos}, \cmd{chaos.sh captura rust-b 10}. Resultado: 5 retries/60~s, \cmd{netem delay 150ms loss 25\%} na saída de B para 10.9.8.10, SYN-ACK em 1,17~s seguido de RST. \textbf{Confirmada} (03:29:11).
  \item \textbf{O monitor parou de atualizar} (evento 4 do enunciado). Teste: \cmd{promq 'time() - watchdog_last_probe_timestamp_seconds'} e \cmd{'up'}; painel. Resultado: idade $<5$~s, \cmd{up=1}, painel atualizado. \textbf{Refutada nesta execução} (o painel passou a mostrar B em vermelho, que é outra coisa).
\end{enumerate}

\textbf{Próxima verificação mais informativa} (escolhida após o primeiro evento): sondar B \emph{a partir de um terceiro ponto} (\cmd{docker compose exec -T toolbox curl -s -m 2 -o /dev/null -w '\%\{http_code\}' http://rust-b:8080/health/ready} e \cmd{/health/live}) e ler \cmd{admin.sh status}. Uma única rodada separa as três primeiras hipóteses: se o terceiro ponto vê 200 e o watchdog vê falha, o problema é do observador ou da rede dele; se vê 503 com \cmd{live} 200, o nó está vivo e \emph{diz} que não está pronto; se nada responde, o nó caiu.

\section{Decisão, riscos e reversibilidade}
\textbf{Decisão (03:29:44):} \emph{retirar B do balanceador} com \cmd{scripts/admin.sh pool drenar rust-b}, mantendo o contêiner vivo para diagnóstico, e \textbf{não reiniciar nada}. Motivos: (a) o impacto no usuário era latência e retries, não erro; (b) reiniciar destruiria o estado que explicava o incidente (flag de prontidão, regra de rede, logs) e não atacaria a causa; (c) a drenagem é desfeita com um comando (\cmd{pool restaurar}) e o script guarda uma cópia do \cmd{nginx.conf} em \cmd{evidencias/}.

\textbf{Riscos aceitos:} A vira ponto único de falha e recebe 100\% do tráfego (3~req/s, muito abaixo da capacidade); se A falhasse durante a drenagem, o serviço cairia. Mitigação: manter o fundo baixo, acompanhar \cmd{status.sh --vivo} e devolver B assim que o watchdog o considerasse UP.

\textbf{O que deu errado e como vimos:} nesta máquina (macOS) o \cmd{sed -i -E} do script é incompatível com o \cmd{sed} BSD; o comando imprimiu um erro, o \cmd{nginx -t} passou (arquivo inalterado) e o \cmd{pool status} ``depois'' saiu igual ao ``antes''. A checagem das 03:30:17 confirmou: B ainda recebia 15\% das requisições, com 9 retries. Como as duas causas imediatas já estavam identificadas, aplicamos a \textbf{correção dirigida} (\cmd{ready rust-b 1}; \cmd{chaos.sh limpar rust-b}), ambas reversíveis e registradas em \cmd{acoes.log}/\cmd{chaos.log}. Depois do incidente corrigimos o script (sed portável, commit ``admin.sh: sed portavel'') e provamos a ação: B drenado \seta{} 20/20 por A; restaurado \seta{} 10/10 (\cmd{hora7-10-teste-drenagem-reversivel.txt}).

\textbf{Como confirmamos a recuperação (03:30:25--03:31:00):} watchdog \cmd{down\seta{}up} para B após duas sondagens boas; nenhum alerta ativo (\cmd{NoDown} resolvido às 03:30:32); 20 requisições alternando A/B com 100\% 200; 30~s de tráfego sem retry e p95 de 1~ms; \cmd{RestartCount=0} nos quatro serviços (nenhum reinício em loop); \cmd{chaos.sh status todos} sem regras; painel verde.

\section{Prevenção}
\begin{itemize}[leftmargin=*]
  \item \textbf{Arquitetura:} o balanceador deve consultar a prontidão. O Nginx de código aberto não sonda \cmd{/health/ready}; propomos um controlador simples (pode ser o próprio watchdog, ou um \emph{sidecar}) que, ao marcar DOWN por \cmd{http_503} \emph{confirmado por um segundo ponto}, aplique \cmd{pool drenar} e, ao voltar a UP por $N$ sondagens, \cmd{pool restaurar}. Alternativa: um balanceador com sondagem ativa (HAProxy \cmd{option httpchk}, Nginx Plus, Envoy).
  \item \textbf{Monitoramento:} medir o caminho do usuário, não só o nó: exportar do log do Nginx (ou de uma sonda sintética que passe pelo Nginx) a latência p95, a taxa de retries (\cmd{upstream} com vírgula) e o status por nó. Foi essa visão que mostrou o segundo problema (rede), invisível em \cmd{sentinel_app_request_duration_seconds}.
  \item \textbf{Sondar de mais de um ponto:} o watchdog e o Nginx viam caminhos diferentes para B; uma segunda sondagem (toolbox ou um segundo watchdog) distingue ``nó diz que não está pronto'' de ``rede até o nó''.
  \item \textbf{Alertas:} (a) \emph{Discordância}: \cmd{watchdog_target_down == 1} e, ao mesmo tempo, requisições 200 do mesmo nó no Nginx por mais de 1~min (sinal de ``vivo, mas não pronto'' ou de observador enganado); (b) \emph{Retries}: taxa de requisições com retry $> 2\%$ por 1~min (hoje só visível em \cmd{status.sh --usuario}); (c) manter \cmd{MonitorParado} com \cmd{absent()} e \cmd{up == 0} (provado na Hora 6); (d) lembrar que \cmd{TaxaDeErroAlta} não enxerga a primeira rajada de uma série nova (Hora 6, exp.~4): inicializar os contadores com zero em cada rota/código no arranque dos nós.
  \item \textbf{Operação:} registrar toda ação administrativa (\cmd{acoes.log}, \cmd{pool.log}, \cmd{chaos.log} já fazem isso) e exigir \cmd{evidence.sh} antes de reiniciar; testar os scripts de resposta no sistema operacional onde serão usados (o erro do \cmd{sed} teria custado minutos num incidente real).
\end{itemize}

\section{Limites da simulação}
Tudo roda em \textbf{um único host, um kernel e um relógio}: não há falha independente de hardware, energia, zona ou roteador, e os horários dos logs são comparáveis só porque todos os contêineres compartilham o relógio. As falhas de rede atuam na \textbf{saída de uma interface virtual} (\cmd{tc netem}), opcionalmente por destino, o que explica por que Nginx e watchdog viram B de formas diferentes; numa rede real, perda e atraso seriam assimétricos, variáveis e também na entrada. O tráfego (3~req/s) não revela saturação, filas nem esgotamento de \emph{workers}. O Docker DNS remove o nome de um contêiner parado (vimos \cmd{reason=dns} em vez de \cmd{timeout}), comportamento que não existe com IPs fixos em máquinas reais. O \cmd{HEALTHCHECK} do Docker roda dentro do contêiner (\cmd{127.0.0.1}), por isso continua \cmd{healthy} em qualquer falha de rede. Este ambiente foi executado em macOS (Docker via OrbStack, kernel Linux 7.2 da VM): o laboratório foi desenhado para Linux, e um dos scripts de resposta (\cmd{admin.sh pool}) falhou por diferença do \cmd{sed}; o mesmo cenário em Linux não teria esse passo.

\textbf{Melhoria priorizada e teste futuro:} prontidão consultada pelo balanceador (controlador watchdog\seta{}\cmd{pool drenar/restaurar}, com confirmação por um segundo ponto), porque foi a lacuna comum ao experimento 6 e ao incidente. Teste futuro: repetir o incidente com o controlador ligado e medir (i) o tempo entre o DOWN e a saída de B do pool (meta $<30$~s), (ii) zero retries vistos pelo usuário após a saída e (iii) retorno automático de B após a correção, sem reinício.

\end{document}
